\documentclass[11pt]{article}
\usepackage[T1]{fontenc}
\usepackage[utf8]{inputenc}
\usepackage{lmodern,amsmath,amssymb,mathtools}
\usepackage[a4paper,margin=27mm]{geometry}
\usepackage{microtype,booktabs,array,fancyhdr}
\usepackage[hidelinks]{hyperref}
\usepackage{url}
\setlength{\headheight}{14pt}
\setlength{\emergencystretch}{2em}
\pagestyle{fancy}\fancyhf{}
\fancyhead[L]{\small Erd\H{o}s \#271: verification supplement}
\fancyhead[R]{\small GitHub archive edition v1.3}
\fancyfoot[C]{\thepage}
\newcommand{\N}{\mathbb N_0}
\newcommand{\RR}{\mathcal R}
\newcommand{\card}[1]{\lvert #1\rvert}
\newcommand{\Cs}{C_*}
\newcommand{\al}{\alpha}
\title{Formal verification and reproducibility supplement}
\author{}
\date{Companion to Paper 1 v2.7 and Paper 2 v1.6; revision 1.3 (GitHub archive edition); 1 November 2026}
\begin{document}\maketitle
This supplement retains the detailed verification and provenance sections from Paper~1 v2.3 and Paper~2 v1.3. The earlier editorial revision changed no Lean source or mathematical theorem. The archived rev~5 and rev~6 evidence below retains its original scope; Section~\ref{sec:rev7-comparator} records the separate rev~7 extension and its current verification status.

\paragraph{What is and is not certified.}
All three main rev-5 results have complete formal proofs and a successful independent project rebuild. The supplementary rigidity limit and exact Dai--Chen comparison are conventional paper proofs, not separate Lean theorems. The stronger $1675\sqrt3/1536$ instance has a separately checked proof, and Cedar's independent rev-6 replay verified its integration in the 39-module extension. The original rev-5 archive remains unchanged. The rev~7 source extension is described separately in Section~\ref{sec:rev7-comparator}.

Correspondence between the conventional theorems and the Lean statements was
established by manual source review, supported by four separately checked
bridge lemmas. Lean checks those lemmas and the formal statements; it does not
machine-check the translation of this prose into Lean.

\paragraph{Citation authenticity.}
The accompanying 15-entry citation audit distinguishes direct primary records
from corroboration in published bibliographies. Thirteen works were matched to
an arXiv, publisher, or original author-archive record. The 1979 Gerver--Ramsey
and 1999 Erd\H{o}s--Lev--Rauzy--S\'andor--S\'ark\"ozy articles are independently
corroborated in several primary papers' reference lists; their original
publisher pages were not accessible in this audit. No fabricated reference was
identified. This is not a claim that every cited proof has been reread or that
literature priority is settled. In particular, Fang's full proof remains an
access-dependent method-comparison task.

\paragraph{Independent test.}
Use the public rev~5 tree at
\url{https://github.com/Dlgt-io/room-math-results/tree/main/results/erdos-271/lean/rev5},
not the earlier bootstrap kit. After checking its manifest, follow the archived
\path{VERIFY.md} at
\url{https://github.com/Dlgt-io/room-math-results/tree/main/results/erdos-271/lean/archived-extras}
and run
\begin{quote}\ttfamily\small
export EXPECT\_SUMS=\textbackslash\\
\path{2d4f22165d1e3b60461315f8cd4e2d8e70d0498aaf1add86e4bd459ff4f090ef}\\
bash lean/rev5/scripts/replay\_fresh.sh /path/to/new-empty-project
\end{quote}
where the target directory must not already exist. The full command with the
expected manifest hash and an optional pinned dependency donor is in the
archive's guide. The acceptance criterion is a successful full build followed
by the strict 24-name axiom audit, regression tests, source-token scan, bundle
composition check, and final source-hash check. A fresh \path{logs/SUCCESS.json}
is written only after all checks pass. Internet access may be needed to install
the pinned toolchain and dependencies; the proof sources do not depend on a
research-room connection.

\section{Paper 1: detailed formalization record}
The main theorem has a complete Lean 4.35.0 formalization, including the standard-seed wrappers. The modular release pins Mathlib at commit
\begin{center}\small\ttfamily 56a005855ea17e0b0958f73ec801ae140b0a1679\end{center}
The following evidence concerns different objects and must not be conflated.

\begin{center}\small
\begin{tabular}{@{}p{0.27\textwidth}p{0.66\textwidth}@{}}
\toprule
Layer & Evidence at this revision\\
\midrule
Modular project & Author build, followed by Cedar's independent fresh-directory rebuild of all 38 rev-5 modules on 22 October 2026. The strict audit passed for all 24 expected declarations, with unchanged source hashes.\\[3pt]
External checker & Seven self-contained proof bundles, B1/B5/B6/B7/B8a/B8b/B9, accepted in the fixed Lean/Mathlib environment. Their actual returned statements and input hashes were independently inspected.\\[3pt]
Bundle composition & All six dependency links match. The exact accepted terms compile together as \path{ComposedC9.composed}, also rebuilt by Cedar. Separately, the checker accepted the implication from the seven bundle statements to $C_9$ (job \texttt{68fdc6fb-0c74-4be9-a243-8d7f6ad400ef}, archived in \path{extras/c9-glue/}; independently resubmitted as job \texttt{e0c9847a}). Together these are eight checker verdicts in one environment joined by modus ponens, not one kernel check: a single-term $C_9$ attempt (job \texttt{bb93cb63}) was rejected for exceeding the default \texttt{maxHeartbeats}.\\[3pt]
Final deductions & Two separately written endgames compile over a shared analytic layer; these are not independent reconstructions of every dependency.\\[3pt]
Statement fidelity & Source reviews cover the covering property, positive versus full counts, indexing and greedy cutoff. Four standalone bridge statements also passed the external checker.\\
\bottomrule
\end{tabular}
\end{center}

The reported modular axiom dependencies are \texttt{propext},
\texttt{Classical.choice}, and \texttt{Quot.sound}. No \texttt{sorry}, added axiom,
\texttt{native\_decide}, or unsafe declaration is used in the reported source
closure. Individual external acceptance does not itself constitute a local full
project build or supply a new transitive axiom report.

For clarity, write $S_i$ for the proposition proved by an unconditional checker
bundle and $C_i$ for the final conclusion of a conditional one. The verified
interfaces are
\[
\begin{gathered}
 B1:S_1,\quad B5:S_5,\quad B6:S_6,\quad B7:S_6\to C_7,\\
 B8a:S_1\to S_5\to C_7\to C_{8a},\quad
 B8b:C_{8a}\to C_8,\quad B9:C_8\to C_9.
\end{gathered}
\]
Their compiled direct application is
\[
 B9\bigl(B8b\bigl(B8a\ B1\ B5\ (B7\ B6)\bigr)\bigr).
\]
Thus no analytic obligation remains unproved in this dependency graph. The
interface check and the compilation of this application are separate checks.
The compiled version uses the actual externally accepted proof terms, not axioms
for their conclusions. Cedar subsequently rebuilt this composed declaration as
part of the full frozen rev-5 project. This should not be confused with submitting
one monolithic proof term to the external checker.

\begin{sloppypar}
The unconditional formal assertion is the explicit witness theorem
\path{Erdos271.IndependentEndgame.arbitrarily_large_witnesses}, equivalently
\path{Erdos271.Final.erdos271_limsup_any} in eventual-envelope form.
The separate real-valued \texttt{Filter.limsup} declaration
\path{Erdos271.Final.erdos271_limsup_ge_any} assumes that the normalized quotient
is bounded above on a tail; if it is not, the extended-real limsup in the paper
is $+\infty$. The infinity of $S$ follows from the covering hypothesis, rather
than being an additional exported declaration.
The standard two-seed consequences are
\path{Erdos271.Stanley.two_seed_limsup} and
\path{Erdos271.Stanley.two_seed_ak}.
The finite identity is
\path{Erdos271.FiniteIntegrals.finite_weighted_integral_identity}; its self-contained
checker bundle B5 passed (job \texttt{2c4ac1f1}). It is built in the modular
source, although it is not a separate name among the 24 pinned headline reports.
A frozen-prefix definition with cutoff $m$ is the usual $S(A_0)$ only when
$m=\max A_0$; the standard wrappers impose this. The additional fidelity lemmas
prove uniqueness of a continuation satisfying the greedy rule, identify the
finite positive count, justify the natural-number subtraction in covering
witnesses, and equate the no-envelope statement with arbitrarily large witnesses.
\end{sloppypar}

\paragraph{From paper lemmas to Lean files.}
The following is a guide to the main reusable statements, not a requirement to inspect the formal proof to understand the paper. The modules are part of the pinned release.
\begin{center}\small
\begin{tabular}{@{}p{0.36\textwidth}p{0.56\textwidth}@{}}
\toprule
Conventional step & Principal Lean source and declaration\\
\midrule
Weak covering, pair count, lower envelope & \path{Audit/WCover.lean}, \path{Audit/Moy.lean}\\
Weighted pair inequality & \path{Audit/MemberTail.lean}; checker bundle B1\\
Finite integral identity (Section 3) & \path{WeightedIdentity.lean}: \path{finite_weighted_integral_identity}; bundle B5\\
Kernel mass, moment and logarithmic window & \path{Audit/KMChecked.lean}, \path{Audit/Glue4K.lean}, \path{Audit/Glue5.lean}, \path{Audit/Glue6.lean}; bundles B6--B7\\
Contradiction and arbitrarily large witnesses & \path{Audit/Final.lean}: \path{limit_of_budget}; \path{Audit/General.lean}: \path{erdos271_limsup_any}; \path{IndependentEndgame.lean}: \path{arbitrarily_large_witnesses}\\
Enumeration and ordinary greedy seeds & \path{Audit/AkForm.lean}, \path{Audit/Stanley.lean}: \path{stanley_std_limsup}, \path{stanley_std_ak}\\
Rigidity's finite inequality & \path{IndependentEndgame.lean}: \path{fixed_error_budget}\\
\bottomrule
\end{tabular}
\end{center}
The externally accepted bundles are composed by type-checked function application, with standard axioms. The frozen project, including that composition, now also has an independent fresh-directory rebuild. The paper's conditions on weak covering, positive versus full counts, natural subtraction, zero-indexed enumeration and the greedy continuation have also been compared to the checked statements. These comparisons are manual scope/fidelity reviews; they are separate from the compilation evidence and are not themselves machine-checked.

The defect-budget inequality retaining $\RR+\kappa D$ is formalized by
\par\noindent\path{Erdos271.IndependentEndgame.fixed_error_budget}.\par
The subsequent
logarithmic-mean limiting argument in Paper~1, Corollary~6.1 and the exact
Dai--Chen comparison are reviewed paper deductions, not separate completed Lean
claims.

\paragraph{Independent rebuild and remaining review.}
Cedar's rev-5 replay completed at 19:46:23 UTC on 22 October 2026, with exit
code zero, 38 freshly compiled project modules, 73 matching payload hashes,
a strict 24-of-24 axiom audit, and 16 passing audit regression tests.
The source manifest is
\par\noindent\path{2d4f22165d1e3b60461315f8cd4e2d8e70d0498aaf1add86e4bd459ff4f090ef}.\par
The project was rebuilt from source using pinned cached dependencies; Lean and
Mathlib themselves were not rebuilt from source. The report and raw evidence are
in the public \path{verification/rev5/} directory at
\url{https://github.com/Dlgt-io/room-math-results/tree/main/results/erdos-271/verification/rev5}.
This evidence concerns rev~5, not rev~6 or a new literature-priority
certificate.

This remains a public review draft, not a submission. Human authorship,
a permanent public archive, a final literature-priority check and external expert
review remain release tasks. The theorem comparison establishes improvement
over Dai--Chen's cited result; absence of a later or unindexed equivalent theorem
has not been exhaustively established.


\section{Paper 2: detailed formalization record}
This is a public review draft, not an externally submitted paper. The following is a status record for the release pinned to Lean 4.35.0 and Mathlib commit \texttt{56a00585}; the full commit hash is given in Section~1 and the release manifest below.

\paragraph{Evidence terminology.}
Here \emph{external check} means a self-contained statement/proof term accepted by the fixed Lean/Mathlib checker. \emph{Modular build} means a theorem compiled in the full pinned source tree in the author's environment. \emph{Independent complete rebuild} means compiling that full source tree in a separate environment. These are distinct assurances and none is asserted without its own evidence.

\paragraph{The limsup theorem.}
The declarations \path{Erdos271.B16.B16_main} and \path{Erdos271.B16.B16_exists} include the recursive set definition, strong covering beyond 88, a uniform counting bound and arbitrarily large lower witnesses for every $U<17/6$. Their self-contained statement/proof pairs were accepted by the fixed remote checker in jobs beginning \texttt{7f84be19} and \texttt{e343345f} (statement SHA-256 prefixes \texttt{cf6d045e} and \texttt{609db29c}). Both checker statements contain the uniform inequality and arbitrary-large-witness formulation, not the \texttt{Filter.limsup} equality. The defining equation is not an unsupported existence assumption: the existence statement is also checked. The modular development additionally contains \path{mem_iff_blocks}, identifying the recursive and block definitions, and \path{B16_limsup}/\path{S16_limsup}, giving the explicit \texttt{Filter.limsup} equality. Those additional declarations were also compiled in Cedar's fresh full-project rebuild; they are not separate standalone checker submissions.

\paragraph{The liminf theorem.}
\path{Erdos271.Liminf.StaticCert.liminf_of_cert} is the generic transfer theorem (Paper~2, Theorem~4.2); its 416-element instance is \path{Erdos271.Liminf.liminf_construction_1889}. These are compiled in the author's pinned modular project and Cedar's independent replay, with standard transitive axiom reports; they are not themselves standalone checker-submitted terms. The external checker separately accepted Claude's complete, self-contained proof of the same $419\sqrt3/384$ existence assertion in job \texttt{d5e58222} (statement SHA-256 prefix \texttt{186d2fca}). It constructs the same finite sets using symbolic lifts and bounded repair checks after deleting 13. Grok independently recompiled that precise proof term and obtained only \texttt{propext}, \texttt{Classical.choice} and \texttt{Quot.sound}. Its slightly looser count admits at most four instead of three future-stage points, which does not change the asymptotic constant. A different direct proof using $S_{768}$ and constant $53\sqrt3/48$ was checker-verified as bundle B4 (job \texttt{0e2507a5}); it is not a remote check of \path{liminf_of_cert}. The newer $1675\sqrt3/1536$ result above has its own separate checker verdict and does not modify the frozen modular project.

\paragraph{Counting conventions.}
The base-16 formal count is $A(x)=\card{A_{16}\cap[1,\lfloor x\rfloor]}$, which omits zero, so $F_{A_{16}}(x)=A(x)+1$. Its inequality $A(x)+1\le(17/6)\sqrt{x}$ is exactly the full-count bound in Paper~2, Theorem~1.1. The locally compiled declarations \path{B16_limsup} and \path{S16_limsup} state $\limsup_{x\to\infty}A(x)/\sqrt x=17/6$ over real $x$; since $F_{A_{16}}(x)/\sqrt x=A(x)/\sqrt x+1/\sqrt x$, they give the same limit superior. In contrast, the formal liminf results use an $\varepsilon$--$N$ statement: one set $B$ is chosen first, and for every $\varepsilon>0$ and every $N\in\mathbb N$ there exists an integer $n\ge N$ with $\card{B\cap[0,n]}\le(c+\varepsilon)\sqrt n$. This implies the real-variable liminf in~Paper~2, equation~(1.3). The count includes any zero; no positivity hypothesis is needed for this implication, although positivity of the constructed set is established internally and in the conventional proof. Every bounded subset of $\N$ is finite, so \texttt{ncard} here equals ordinary finite cardinality.

\paragraph{What has and has not been rebuilt.}
The frozen 38-module release has a clean build and fresh-directory replay in the author's environment, pinned by the release \path{SHA256SUMS} file (SHA-256 \texttt{2d4f22165d1e3b60461315f8cd4e2d8e70d0498aaf1add86e4bd459ff4f090ef}). Only hashed source files were transferred into the fresh directory; no locally compiled project artifacts were reused. Its strict audit accepts exactly 24 named headline theorems, each appearing once, with only \texttt{propext}, \texttt{Classical.choice} and \texttt{Quot.sound}. An audit test suite includes duplicate/missing-name and unexpected-axiom rejection. Cedar completed the separate fresh-directory rebuild at 19:46:23 UTC on 22 October 2026. All 38 project modules were compiled, all 73 payload hashes matched before and after the run, the strict audit passed for 24 distinct expected names, and all 16 regression tests passed. The runner exited with code zero and produced \path{logs/SUCCESS.json}. It used pinned cached dependencies: Lean and Mathlib were not rebuilt from source. The report and raw logs are available under \path{verification/rev5/} in the public archive. This independently reproduced build concerns the frozen rev~5 project; it does not automatically certify the rev~6 additions, literature priority, or a fresh informal-to-formal review.

\paragraph{Navigating the Lean development.}
For a human reader, the recommended order is the main theorem, its two or three central lemmas, and only then the finite certificates. In the pinned release the corresponding statements are:
\begin{center}\small
\begin{tabular}{@{}p{0.34\textwidth}p{0.58\textwidth}@{}}
\toprule
Paper step & Lean file / theorem\\
\midrule
Base-16 definition and covering & \path{B16.lean}: \path{S16_char}, \path{mem_iff_blocks}, \path{S16_cover}, \path{B16_main}, \path{B16_exists}\\
Digit count and global estimate & \path{B16.lean}: \path{N_block}, \path{N_16m}, \path{encode}, \path{member_bound}, \path{bound_all}\\
Limsup equality & \path{B16Limsup.lean}: \path{B16_limsup}, \path{S16_limsup}\\
Finite lift and repair & \path{LiminfCore.lean}: \path{lift_cov}, \path{step_cov}, \path{iter_spec}\\
Generic finite-to-infinite theorem & \path{LiminfGen.lean}: \path{StaticCert.liminf_of_cert}\\
Frozen 416-element certificate & \path{LiminfCert49152.lean}: \path{S49152_cov}, \path{S49152_card}\\
Frozen liminf constant & \path{LiminfFinal.lean}: \path{liminf_construction_1889}\\
\bottomrule
\end{tabular}
\end{center}
The manuscript proofs are independent of any familiarity with Lean. Finite witness lists and other generated data are isolated from the conceptual lemmas. A standalone checker term may be much longer than the corresponding modular proof because the external service does not accept project imports; that should not be mistaken for the intended reader-facing source architecture.

\section{Contribution and literature-access record}
The base-16 construction and finite static certificates originate in Claude's room notes; the complemented-digit counting argument was developed by ChatGPT; Grok assembled the concrete Lean constructions and the generic lift-and-stack formalization. These are contribution records for the collaboration, not a settled publication author list.

This draft freezes the $17/6$ and $419\sqrt3/384$ targets. Earlier constructions are not advertised as additional release theorems. The smaller $1675\sqrt3/1536$ instance is now separately checker-verified, but lies outside the frozen 38-module source and is presented only as a clearly labeled remark. The independent frozen-project rebuild is complete. Before submission, the author list, final literature-priority review, permanent public source release and external specialist review remain to be settled. The comparison with Fang uses his publisher's precise theorem statement; his full subscription article was not available to the manuscript assembler, so the comparison of proof methods remains open. The full Kiss--S\'andor--Yang and Chen preprints were consulted for definitions and comparison statements. A 2026 preprint of Saengrungkongka (arXiv:2609.02041) was reviewed by a collaborator; its main theorems concern $k\ge6$ or higher-order covering/Kakeya questions, and no overlap with our three stated constants was found. This is not an exhaustive priority check for the current $k=3$ bounds.


\section{Why unordered integer saturation is not forward covering}
This comparison concerns Csajb\'ok--Nagy's integer construction
(Theorem~4.13 of \path{arXiv:2401.06283v4}). Its finite-group results are not
needed here. Let $H_l$ contain the length-$l$ base-4 integers with digits in
$\{2,3\}$, for $l\ge1$, and put $S=\bigcup_{l\ge1}H_l$. Define
\[
x_l=\frac{2(4^{l+1}-1)}3-1,\qquad
N_l=|S\cap[1,x_l]|=2^{l+1}-2.
\]
There are $x_l-N_l$ non-members in the prefix. A forward witness
$a<b<n\le x_l$ uses a pair from that same prefix, and each pair represents
only one target. Thus the number of non-members without forward witnesses is
at least
\[
x_l-N_l-\binom{N_l}{2}
=\frac{2\cdot4^l+9\cdot2^l-8}{3}.
\]
Dividing by $x_l$ gives a lower bound tending to $1/4$. This occurs along
these cutoffs; it is not an assertion about every individual digit block.
Deleting a fixed initial interval cannot remove all these failures. Hence the
unordered saturation property of this set does not give eventual weak forward
covering. The argument is a conventional counting proof, not a new Lean result
or a claim of novel saturation theory.

\section{Identity of the separately checked stronger instance}
The full checker job for $1675\sqrt3/1536$ is
\par\noindent\path{042a1cf7-e251-4071-90f5-b1f384498329}.\par
Its statement SHA-256 is
\par\noindent\path{f34ea2a577ba67182a137abd724314a1680934c5b5e52e9886a1145a90589836}\par
and its proof SHA-256 is
\par\noindent\path{4cab2175f543161e3fff9819fb22bf70aa39920232f83be2bfbec39009ef81b3}.\par
The adopted rev-6 manifest is
\par\noindent\path{78291fa4ca02d51370c4b2c56073ef58bb5f98d7553a7fa77c5fc156d2219c03}.\par
Cedar's independent replay completed at 12:42:29 UTC on 23 October 2026,
with exit code zero, 77 matching payload hashes, 39 freshly compiled project
modules, 26 distinct audited declarations, and all 17 audit regression tests
passed with the real transcript. The printed types and axiom reports for the
original 24 declarations were byte-identical to the accepted rev-5 run.
The receipts are \path{verification/rev6/SUCCESS.json} and
\path{verification/comparison-and-acceptance.json} in the public archive at
\url{https://github.com/Dlgt-io/room-math-results/tree/main/results/erdos-271/verification}.
The October 2026 archive edition
inspected those receipts; that editorial round did not perform another Lean build.
The replay used pinned cached dependencies, not a Lean/Mathlib source rebuild.
Rev~5 and its successful verification remain unchanged.
\section{Rev~7 challenge and solution comparator}\label{sec:rev7-comparator}
Rev~7 is a GitHub-only extension of the weighted-covering formalization. It preserves the rev~6 manuscript and Lean sources and introduces no new mathematical bound or research-room adoption. \path{Statement.lean} defines the inclusive counting function $F_S(x)=|S\cap[0,x]|$, returning zero when $x<0$, and defines weak $AP_3$-covering by explicit witnesses $a<b<z$ with $z=2b-a$. \path{Challenge.lean} imports \path{Statement.lean} alone; the statement module depends on Mathlib, not on the solution or legacy proof modules.

The challenge theorem is \path{Erdos271.main_result}. It asserts the abstract's lower bound $\sqrt{2\sqrt2/\log(1+\sqrt2)}$ for the extended-real limsup of $F_S(x)/\sqrt{x}$ under weak covering, with no eventual-boundedness assumption and no \texttt{DecidablePred} hypothesis. \path{Solution.lean} imports \path{Abstract.lean} and the legacy proof modules to prove this theorem. \path{Abstract.lean} supplies the bridge from the legacy positive count and real-valued endgame to the inclusive count and extended-real conclusion; the theorem's mathematical content is unchanged.

\paragraph{Comparator protocol and status.}
On 1 November 2026, native \texttt{lake comparator} from Lean 4.35.0 accepted \path{Erdos271.main_result}. It compared the trusted challenge with the solution, checked the permitted axioms, and replayed the exported proof using both Lean's kernel and the independent \texttt{nanoda} kernel. The configuration permits only \texttt{propext}, \texttt{Quot.sound}, and \texttt{Classical.choice}. Controls accepted a valid proof and rejected a changed theorem statement and an extra axiom.

The checking environment was prepared from a clean Linux runtime and a fresh official Lean release. Pinned Mathlib dependencies and their public cache were fetched afresh. Comparator ran before any ordinary proof build; no project build outputs were copied from an earlier run. Later project build and audit commands run in a nonnetwork bubblewrap sandbox. Project sources, checker binaries, and parent-owned evidence are read-only, while the project's \path{.lake} build directory and temporary scratch space are writable. The sandbox controls check non-root execution, zero effective and permitted capabilities, no-new-privileges, read-only evidence, and absence of an external network route.

The scripts record raw comparator output, standard error, exit code, tool versions, and source and checker hashes under \path{lean/rev7/logs/comparator/}, outside the proof-writable \path{.lake} directory. The archived evidence and combined replay receipt are under \path{verification/rev7/}. The complete replay also builds the inherited constructions and audits 27 designated declarations. The earlier rev~6 replay and standalone checker receipts retain their historical scope. Cached dependencies were used; Lean and Mathlib were not rebuilt from source. The trusted statement and referenced definitions still require mathematical review.
\end{document}
